\documentclass[11pt]{article}
\usepackage{reusv}

\setborderthickness{0.5pt}
\setbordermargin{1cm}
\setbordercolor{black}

\slimtitle{베지어 제약조건과 컨벡스화}

\begin{document}

\drawborder
\thispagestyle{firstpage}

\lightertitle{베지어 제약조건과 컨벡스화}

{\normalsize\textit{Research Note No.~005 $|$ bezier-orbit-transfer-scp}}\par\medskip

\def\lighterdate{March 12, 2026}
\lighterauthor{Suwon Lee$^{\dagger}$}

\lighteraddress{$^\dagger$}{Future Mobility Control Lab, Kookmin University, Seoul, Korea}
\lighteremail{suwon.lee@kookmin.ac.kr}

\begin{abstract}
본 보고서는 베지어 곡선으로 성형된 추진 가속도에 부과되는 제약조건을
컨벡스 최적화 프레임워크 안에서 처리하는 방법을 정리한다.
볼록껍질 성질(convex hull property)을 이용한 보수적 제약 처리와,
극값 정리(extremum theorem)를 이용한 무손실 컨벡스화(lossless convexification)를
비교하여 제시한다.
또한 추력 크기 제약 $\|\mathbf{u}(\tau)\| \leq u_{\max}$를
2차 원뿔 계획법(SOCP) 제약으로 변환하는 절차를 유도한다.
\end{abstract}

%% ============================================================
\section{문제 설정}
%% ============================================================

\cite{report003}에서 채택한 추진 가속도 베지어 성형에서
설계변수는 제어점 행렬 $\mathbf{P}_u \in \mathbb{R}^{(N+1)\times 3}$이다.
물리적 제약조건은 연속 시간 $\tau \in [0,1]$ 전체에 걸쳐 만족되어야 하므로,
유한 개의 최적화 제약조건으로 변환하는 것이 필요하다.

다루는 주요 제약조건은 두 종류이다.

\begin{description}
  \item[추력 크기 제약] $\|\mathbf{u}(\tau)\| \leq u_{\max}$,\quad $\forall\tau\in[0,1]$
  \item[축별 부등식 제약] $u_k(\tau) \leq c_k$,\quad $\forall\tau\in[0,1]$,\quad $k \in \{x,y,z\}$
\end{description}

%% ============================================================
\section{볼록껍질 성질을 이용한 보수적 제약}
%% ============================================================

\subsection{볼록껍질 성질}

Bernstein 기저의 분할 단위 성질과 비음성으로부터,
베지어 곡선은 제어점들의 볼록껍질(convex hull) 안에 항상 존재한다.

\begin{equation}
  q(\tau) \in \mathrm{conv}\{p_0, p_1, \ldots, p_N\}, \quad \forall\tau\in[0,1]
  \label{eq:convex_hull}
\end{equation}

\subsection{보수적 제약 처리}

스칼라 상한 제약 $q(\tau) \leq c$에 대해 볼록껍질 성질을 적용하면:

\begin{equation}
  p_i \leq c,\quad \forall\, i = 0, 1, \ldots, N
  \quad \Rightarrow \quad q(\tau) \leq c,\quad \forall\tau\in[0,1]
  \label{eq:conservative}
\end{equation}

모든 제어점이 상한 이하면 곡선 전체도 상한 이하가 됨이 보장된다.
이 조건은 제어점에 대한 선형 부등식이므로 LP/QP에 직접 포함할 수 있다.

단, 이 조건은 충분조건이지 필요충분조건이 아니다.
제어점이 경계를 넘더라도 곡선 자체는 경계 안에 있을 수 있으므로,
볼록껍질 조건은 실제 허용 영역보다 작은 영역만 탐색하는
\textbf{보수적(conservative)} 접근이다.

%% ============================================================
\section{극값 정리를 이용한 무손실 컨벡스화}
%% ============================================================

\subsection{극값 정리 적용}

연속 함수 $q(\tau)$가 $[0,1]$에서 상한 $c$를 넘지 않을 조건은
다음과 동치이다.

\begin{equation}
  q(\tau) \leq c,\;\forall\tau\in[0,1]
  \;\Longleftrightarrow\;
  q(\tau^*) \leq c,\;\forall\tau^* \in \mathcal{T}^*
  \label{eq:extremum}
\end{equation}

여기서 $\mathcal{T}^* = \{0,\,1\} \cup \{\tau : q'(\tau) = 0\}$는
끝점과 내부 극값점의 집합이다.
즉, 연속 제약을 극값점에서의 유한 개 제약으로 정확히 대체할 수 있다.

\subsection{극값점 계산}

$q'(\tau) = \mathbf{B}_{N-1}(\tau)^T (\mathbf{D}_N \mathbf{p}) = 0$의 근을 구해야 한다.
$\mathbf{D}_N \mathbf{p} \in \mathbb{R}^N$을 계수로 하는 차수 $N-1$ 다항식의 근이므로,
\textbf{동반 행렬(companion matrix)}의 고유값으로 안정적으로 계산할 수 있다.
실수 근이면서 $\tau^* \in (0,1)$인 것만 $\mathcal{T}^*$에 포함한다.

\subsection{SCP에서의 적용}

SCP 반복에서 참조 제어점 $\bar{\mathbf{p}}$가 주어지면:
\begin{enumerate}
  \item $\mathbf{D}_N \bar{\mathbf{p}}$의 근 $\{\bar{\tau}_j^*\}$를 계산한다.
  \item $\mathcal{T}^* = \{0, 1, \bar{\tau}_1^*, \ldots, \bar{\tau}_m^*\}$에서
        Bernstein 기저를 평가한다.
  \item 제약행렬 $\mathbf{C} = [\mathbf{B}_N(\bar{\tau}^*_j)^T]_j \in \mathbb{R}^{|\mathcal{T}^*|\times(N+1)}$을 구성한다.
  \item 선형 제약 $\mathbf{C}\,\mathbf{p} \leq c\,\mathbf{1}$을 QP/SOCP에 추가한다.
\end{enumerate}

이 방법은 참조 궤도에서의 극값점을 기준으로 구성하므로 SCP와 자연스럽게 결합되며,
볼록껍질 조건과 달리 보수성이 없다.

%% ============================================================
\section{추력 크기 제약의 SOCP 변환}
%% ============================================================

\subsection{연속 제약의 이산화 필요성}

추력 크기 제약은 노름 제약이므로 각 축 제약의 단순 조합으로 표현되지 않는다.

\begin{equation}
  \|\mathbf{u}(\tau)\|^2 = u_x(\tau)^2 + u_y(\tau)^2 + u_z(\tau)^2 \leq u_{\max}^2
  \label{eq:thrust_norm}
\end{equation}

이를 $\tau \in [0,1]$에서 연속적으로 만족하려면, 극값 정리를 3차원으로 확장하거나
충분히 촘촘한 이산 격자점 $\{\tau_j\}$에서 제약을 부과하는 방법을 쓴다.

\subsection{이산 격자에서의 SOCP 제약}

이산화된 각 시점 $\tau_j$에서 추력 크기 제약은 2차 원뿔 제약이다.

\begin{equation}
  \|\mathbf{u}(\tau_j)\| \leq u_{\max}
  \;\Longleftrightarrow\;
  \left\|\mathbf{P}_u^T \mathbf{B}_N(\tau_j)\right\| \leq u_{\max}
  \label{eq:socp_constraint}
\end{equation}

이는 제어점 행렬 $\mathbf{P}_u$에 대한 2차 원뿔 제약이므로 SOCP로 처리된다.
격자 간격이 충분히 작으면 연속 제약에 대한 근사 정확도가 보장된다.
실용적으로는 $\tau_j = j/M$, $j = 0, 1, \ldots, M$으로 균등 이산화하며
$M = 2N \sim 4N$ 정도가 적절하다.

\subsection{추력 하한 제약 (최소 추력)}

추력 하한 $\|\mathbf{u}(\tau)\| \geq u_{\min}$은 비볼록 제약이다.
Lossless convexification 기법~\cite{report007}에서 이를 등식 제약으로 완화하거나
SCP 외부 루프에서 검증하는 방식을 사용할 수 있다.
본 프로젝트 기본 설정에서는 추력 하한을 고려하지 않는다.

%% ============================================================
\section{두 방법의 비교}
%% ============================================================

\begin{table}[h]
  \centering
  \caption{제약조건 처리 방법 비교}
  \label{tab:constraint_comparison}
  \begin{tabular}{p{3.0cm}p{4.5cm}p{4.5cm}}
    \toprule
    & 볼록껍질 (보수적) & 극값 정리 (무손실) \\
    \midrule
    제약 형태    & 선형 ($N+1$개) & 선형 (극값점 수만큼) \\
    보수성       & 있음 & 없음 \\
    계산 비용    & 없음 (사전 고정) & 극값점 계산 필요 \\
    SCP 의존성   & 독립 & 참조 궤도에 의존 \\
    추천 용도    & 초기 guess, 빠른 계산 & 최종 해 품질 향상 \\
    \bottomrule
  \end{tabular}
\end{table}

%% ============================================================
\section{요약}
%% ============================================================

\begin{itemize}
  \item 볼록껍질 성질은 제어점의 선형 부등식으로 연속 제약을 보수적으로 처리한다.
  \item 극값 정리는 연속 제약을 극값점에서의 유한 조건으로 정확히 대체하여
        보수성 없이 컨벡스 처리를 가능하게 한다.
  \item 추력 크기 제약은 이산화된 격자점에서 SOCP 제약으로 처리한다.
  \item SCP 반복마다 참조 궤도 기준으로 극값점을 갱신하여 정확도를 높인다.
  \item 구체적인 SCP 알고리즘 구조 안에서 이 제약들이 어떻게 조합되는지는
        \cite{report006}에서 다룬다.
\end{itemize}

\bibliographystyle{IEEEtran}
\bibliography{references}

\end{document}
